% 現行発展版（2026-10-15、68cb7492）の組版設定を継承。
% jsbook の 11pt で用いられる拡大率 (\mag) が用紙にも及ぶのを防ぎ、
% PDF の実寸を A4 (210 mm × 297 mm) に保つ。
\documentclass[
  dvipdfmx,
  a4paper,
  10pt,
  nomag,
  fleqn,
  openany
]{jsbook}

% =========================================================
% 数学・作図
% =========================================================

\usepackage{amsmath,amsthm,amssymb,amsfonts}
% BOOK-FILE-ID: f-728762483022
% 基本編・発展編・新解釈で共通の数式設定。
% amsmath の読み込み後に input する。
\everymath{\displaystyle}

% align* は本文を先読みするため、終了マクロまでの本文をまとめて渡す。
% 単純な \begin / \end の別名では amsmath が終了位置を検出できない。
\long\def\balign#1\ealigin{\begin{align*}#1\end{align*}}
% 終了トークンは上のマクロに取り込まれる。単独使用は誤記として検出する。
\newcommand{\ealigin}{\PackageError{book-math}
  {Unmatched \string\ealigin}{Use \string\balign ... \string\ealigin.}}

% enumerate の任意引数（enumitem の設定）もそのまま使用できる。
\newcommand{\benumrate}{\begin{enumerate}}
\newcommand{\eenumrate}{\end{enumerate}}
\newcommand{\bcases}{\begin{cases}}
\newcommand{\ecases}{\end{cases}}


\usepackage{refcount}
\usepackage{tikz}
\usepackage{graphicx}
% 現行発展版と同じ本文内配置・キャプション形式。
\usetikzlibrary{arrows.meta,calc,positioning}
\newcommand{\BookFigureCaption}[1]{%
  \refstepcounter{figure}\par\smallskip
  {\footnotesize 図\thefigure\quad #1}\par}
\newcommand{\BookGraph}[2]{%
  \par\smallskip\noindent
  \begin{minipage}{\linewidth}\centering
    \begin{tikzpicture}[scale=0.8,>=Stealth,line cap=round,line join=round,
      every node/.style={font=\small}]#2\end{tikzpicture}
    \BookFigureCaption{#1}
  \end{minipage}\par\smallskip}
\newcommand{\NewFigBasis}{\BookGraph{同じ点を斜めの基底で記録する}{
 \draw[->](-.6,0)--(3,0)node[right]{$x$};
 \draw[->](0,-1.3)--(0,2.3)node[above]{$y$};
 \draw[->,thick](0,0)--(1,1)node[above left]{$\boldsymbol u$};
 \draw[->,thick](0,0)--(1,-1)node[below]{$\boldsymbol w$};
 \draw[->,very thick](0,0)--(2,1)node[right]{$(2,1)$};
 \draw[dashed](1.5,1.5)--(2,1);
 \draw[dashed](0,0)--(1.5,1.5);
 \node[above]at(1.5,1.5){$\frac32\boldsymbol u$};
 \node[below left]at(0,0){$O$};
}}
\newcommand{\NewFigFixed}{\BookGraph{原点を不動点へ移す}{
 \draw[->](-1.6,1.3)--(3.7,1.3);
 \fill(-.8,1.3)circle(1.5pt)node[below]{$\alpha=-1$};
 \fill(1.8,1.3)circle(1.5pt)node[below]{$a=1$};
 \draw[<->](-.8,1.85)--(1.8,1.85)node[midway,above]{ずれは2};
 \draw[->](-1.6,-.3)--(3.7,-.3);
 \fill(-.8,-.3)circle(1.5pt)node[below]{$0$};
 \fill(1.8,-.3)circle(1.5pt)node[below]{$b=a-\alpha=2$};
 \draw[->,dashed](-.8,.85)--(-.8,.1);
}}
\newcommand{\NewFigRotation}{\BookGraph{複素数の積は回転と拡大を表す}{
 \draw[->](-1.6,0)--(2,0)node[right]{$\operatorname{Re}$};
 \draw[->](0,-1.4)--(0,2)node[above]{$\operatorname{Im}$};
 \draw[gray](0,0)circle(1);
 \draw[gray,dashed](0,0)circle(1.4);
 \draw[->,thick](0,0)--(25:1)node[right]{$z$};
 \draw[->,very thick](0,0)--(75:1.4)node[above right]{$\lambda z$};
 \draw[->](25:.65)arc(25:75:.65);
 \node at(50:.92){$\theta$};
 \node[below left]at(0,0){$O$};
}}
\newcommand{\NewFigRatio}{\BookGraph{二つの基準点からの差を比べる}{
 \coordinate(a)at(-1.2,0);
 \coordinate(b)at(1.4,-.4);
 \coordinate(z)at(.3,1.8);
 \draw[->,thick](a)--(z)node[midway,left]{$z-\alpha$};
 \draw[->,thick](b)--(z)node[midway,right]{$z-\beta$};
 \draw[dashed](a)--(b);
 \fill(a)circle(1.5pt)node[below]{$\alpha$};
 \fill(b)circle(1.5pt)node[below]{$\beta$};
 \fill(z)circle(1.5pt)node[above]{$z$};
}}
\newcommand{\NewFigCayley}{\BookGraph{傾きから単位円の点へ}{
 \draw[->](-1.5,0)--(1.6,0)node[right]{$u$};
 \draw[->](0,-1.3)--(0,1.5)node[above]{$v$};
 \draw(0,0)circle(1);
 \draw[thick](-1,0)--(.5,.866);
 \fill(-1,0)circle(1.5pt)node[below left]{$-1$};
 \fill(.5,.866)circle(1.5pt)node[above right]{$\Phi(a)$};
 \fill(0,.577)circle(1.2pt)node[left]{$(0,a)$};
 \draw[->](-.55,0)arc(0:30:.45);
 \node at(-.40,.17){$\theta$};
 \draw[->](.38,0)arc(0:60:.38);
 \node at(.61,.30){$2\theta$};
}}
\newcommand{\NewFigSpiral}{\BookGraph{一定倍率と一定角で進む状態}{
 \begin{scope}[scale=.52]
 \draw[->](-4.8,0)--(3,0)node[right]{$\operatorname{Re}$};
 \draw[->](0,-4.8)--(0,3.2)node[above]{$\operatorname{Im}$};
 \draw[->,thick](1,0)--(1,1);
 \draw[->,thick](1,1)--(0,2);
 \draw[->,thick](0,2)--(-2,2);
 \draw[->,thick](-2,2)--(-4,0);
 \draw[->,thick](-4,0)--(-4,-4);
 \foreach \x/\y in {1/0,1/1,0/2,-2/2,-4/0,-4/-4}
  {\fill(\x,\y)circle(2pt);}
 \node[below right]at(1,0){$z_1$};
 \node[right]at(1,1){$z_2$};
 \node[above right]at(0,2){$z_3$};
 \node[above]at(-2,2){$z_4$};
 \end{scope}
}}
\newcommand{\NewFigDoubling}{\BookGraph{角度の倍化は一定角の回転と異なる}{
 \draw[->](-1.4,0)--(1.6,0);
 \draw[->](0,-1.2)--(0,1.5);
 \draw(0,0)circle(1);
 \draw[thick](0,0)--(30:1);
 \draw[thick](0,0)--(60:1);
 \draw[thick](0,0)--(120:1);
 \fill(30:1)circle(1.5pt)node[right]{$z$};
 \fill(60:1)circle(1.5pt)node[above right]{$z^2$};
 \fill(120:1)circle(1.5pt)node[above left]{$z^4$};
 \draw[->](30:1.18)arc(30:60:1.18);
 \draw[->](60:1.3)arc(60:120:1.3);
}}
\newcommand{\NewFigProjective}{\BookGraph{長さが変わっても成分比は変わらない}{
 \draw[->](-.4,0)--(3.7,0)node[right]{$x$};
 \draw[->](0,-.3)--(0,2.6)node[above]{$y$};
 \draw[dashed](0,0)--(3.3,2.2);
 \draw[->,thick](0,0)--(1.5,1)node[above left]{$(x,y)$};
 \draw[->,very thick](0,0)--(3,2)node[above left]{$(cx,cy)$};
 \node[below right]at(1.5,.3){$\frac{x}{y}=\frac{cx}{cy}$};
}}



% =========================================================
% レイアウト
% =========================================================

\ifdefined\BookMobile
  \usepackage[paperwidth=90truemm,paperheight=160truemm,margin=8truemm,headheight=8.5truemm,headsep=1truemm]{geometry}
\else
  \usepackage[margin=20truemm]{geometry}
\fi

\usepackage{multicol}
% 本文は同じ原稿を使用し、スマホ版だけ1段組にする。
\ifdefined\BookMobile
  \newenvironment{bookcolumns}{\par}{\par}
\else
  \newenvironment{bookcolumns}{\begin{multicols*}{2}}{\end{multicols*}}
\fi
\usepackage{enumitem}
\usepackage{etoolbox}
\usepackage{url}

% =========================================================
% 見出し・章デザイン
% =========================================================

\usepackage{quotchap}


% =========================================================
% 色・ボックス
% =========================================================

\usepackage[table]{xcolor}
\usepackage{tcolorbox}

\tcbuselibrary{breakable}

% 発想・条件・結論を短く読み返すための左罫線。本文の証明は外に置く。
\newtcolorbox{bookpoint}{unbreakable,sharp corners,
  colback=black!3,colframe=black!55,boxrule=0pt,leftrule=1.2pt,
  left=5pt,right=3pt,top=3pt,bottom=3pt,
  before skip=6pt,after skip=6pt}
\newcommand{\bookstep}[1]{\par\medskip\noindent\textbf{#1}\par\nobreak\smallskip}

\definecolor{shadecolor}{gray}{0.80}


% =========================================================
% 全体の組版
% =========================================================

% ページの高さを無理にそろえず、上から順に流す
\raggedbottom
\raggedcolumns

% 参照用の段組を使う場合の段間
\setlength{\columnsep}{12pt}

% 段組の中央線
\setlength{\columnseprule}{0.4pt}

% 段落間隔
\setlength{\parskip}{0.35\baselineskip plus 1pt}

% 行送りの不自然な拡張を必要な箇所だけ許容
\setlength{\emergencystretch}{1em}
\linespread{1.05}


% =========================================================
% 数式周辺の余白
% =========================================================

\setlength{\abovedisplayskip}{10pt plus 2pt minus 2pt}
\setlength{\belowdisplayskip}{10pt plus 2pt minus 2pt}

\setlength{\abovedisplayshortskip}{8pt plus 2pt minus 2pt}
\setlength{\belowdisplayshortskip}{8pt plus 2pt minus 2pt}


% =========================================================
% 箇条書き
% =========================================================

% 例題・方針・解答・解法などで使用している
% enumerate のインデントは元の設計を維持する。
%
% leftmargin を指定しないことで、jsbook / enumitem の
% \leftmargini を基準とした元来の狭い幅を保持する。
\setlist[enumerate,1]{
  label=(\arabic*),
  leftmargin=\leftmargini,
  parsep=3pt,
  beginpenalty=10000
}


% =========================================================
% 引用
% =========================================================

% enumerate のインデント体系と整合するよう、
% \leftmargini を基準に smallquote の幅を決める。
\newlength{\smallquotemargin}

\setlength{\smallquotemargin}{0.5\leftmargini}

\makeatletter
\AtBeginEnvironment{quote}{\@beginparpenalty=10000}
\newenvironment{smallquote}
  {%
    \begin{list}{}{%
      \@beginparpenalty=10000%
      \setlength{\leftmargin}{\smallquotemargin}%
      \setlength{\rightmargin}{\smallquotemargin}%
      \setlength{\topsep}{0.35\baselineskip}%
      \setlength{\partopsep}{0pt}%
      \setlength{\parsep}{0pt}%
      \setlength{\listparindent}{0pt}%
    }%
    \item[]%
  }
  {\end{list}}
\makeatother


% =========================================================
% 箇条書き内で式番号を付ける
% =========================================================

\newcommand{\itemeq}[2][]{%
  \refstepcounter{equation}%
  \item
  \makebox[\linewidth]{%
    $\displaystyle #2$
    \hfill
    (\theequation)%
  }%
  \ifx\relax#1\relax
  \else
    \label{#1}%
  \fi
}


% =========================================================
% 証明終了記号
% =========================================================

\renewcommand{\qed}{%
  \unskip\nobreak\quad\qedsymbol
}


% =========================================================
% 目次・見出し番号
% =========================================================

\setcounter{tocdepth}{3}
\setcounter{secnumdepth}{4}

% 両版共通のPDF内リンク。パッケージは他の設定の後で読み込む。
\usepackage[dvipdfmx,bookmarksnumbered,bookmarksopen,linktoc=all,hidelinks]{hyperref}
\usepackage{pxjahyper}
\usepackage[normalem]{ulem}

% 文字色と字幅を変えず、移動の入口だけに控えめな点線を付ける。
% 数式中では下線を描かない。
\DeclareRobustCommand{\booklinkcue}[1]{%
  \bookcueprepared{#1}%
}
\DeclareRobustCommand{\booknavref}[2]{\hyperref[#1]{\booklinkcue{#2}}}

% 用語へのリンクはすべて残す。各節・項で最初に本文に現れる語だけを装飾する。
% 見出し、目次、見出し参照での出現は本文の初出として数えない。
\newif\ifbooktermcues
\booktermcuestrue
\ExplSyntaxOn
\tl_new:N \l_book_cue_text_tl
\int_new:N \l_book_cue_previous_int
\int_new:N \l_book_cue_current_int
\int_new:N \l_book_cue_penalty_int
\cs_new_protected:Npn \bookcueprepared #1
  {
    \ifmmode #1\else
      \bgroup
      \tl_clear:N \l_book_cue_text_tl
      \int_zero:N \l_book_cue_previous_int
      \tl_map_inline:nn {#1}
        {
          \token_if_cs:NTF ##1
            { \int_zero:N \l_book_cue_previous_int }
            {
              \int_set:Nn \l_book_cue_current_int { `##1 }
              \int_compare:nNnT {\l_book_cue_current_int} > {255}
                {
                  \int_compare:nNnT {\l_book_cue_previous_int} > {255}
                    {
                      % ulemの語の中でも、日本語の禁則に従って改行する。
                      \int_set:Nn \l_book_cue_penalty_int
                        {
                          \tex_prebreakpenalty:D \l_book_cue_current_int
                          + \tex_postbreakpenalty:D \l_book_cue_previous_int
                        }
                      \tl_put_right:Nx \l_book_cue_text_tl
                        { \exp_not:N \penalty \int_use:N \l_book_cue_penalty_int \exp_not:N \relax }
                    }
                }
              \int_set_eq:NN \l_book_cue_previous_int \l_book_cue_current_int
            }
          \tl_put_right:Nn \l_book_cue_text_tl {##1}
        }
      \markoverwith{\textcolor{black!45}{\rule[-0.7ex]{0.3pt}{0.3pt}}\kern1.5pt}
      \exp_args:NV \ULon \l_book_cue_text_tl
    \fi
  }
\prop_new:N \g_book_seen_terms_prop
\tl_new:N \l_book_term_key_tl
\DeclareRobustCommand{\bookterm}[2]
  {
    \ifmmode
      \hyperlink{term:#1}{#2}
    \else
      \ifbooktermcues
        \tl_set:Nx \l_book_term_key_tl
          { \thechapter / \thesection / \thesubsection / #1 }
        \exp_args:NNV \prop_if_in:NnTF \g_book_seen_terms_prop \l_book_term_key_tl
          { \hyperlink{term:#1}{#2} }
          {
            \exp_args:NNV \prop_gput:Nnn \g_book_seen_terms_prop \l_book_term_key_tl {seen}
            \hyperlink{term:#1}{\booklinkcue{#2}}
          }
      \else
        \hyperlink{term:#1}{#2}
      \fi
    \fi
  }
\ExplSyntaxOff
\makeatletter
\pretocmd{\@sect}{\booktermcuesfalse}{}{\PackageError{book-links}{Cannot suppress heading cues}{}}
\apptocmd{\@sect}{\booktermcuestrue}{}{\PackageError{book-links}{Cannot restore body cues}{}}
\pretocmd{\@ssect}{\booktermcuesfalse}{}{\PackageError{book-links}{Cannot suppress starred heading cues}{}}
\apptocmd{\@ssect}{\booktermcuestrue}{}{\PackageError{book-links}{Cannot restore starred body cues}{}}
\pretocmd{\@makechapterhead}{\booktermcuesfalse}{}{\PackageError{book-links}{Cannot suppress chapter cues}{}}
\apptocmd{\@makechapterhead}{\booktermcuestrue}{}{\PackageError{book-links}{Cannot restore chapter body cues}{}}
\pretocmd{\@makeschapterhead}{\booktermcuesfalse}{}{\PackageError{book-links}{Cannot suppress starred chapter cues}{}}
\apptocmd{\@makeschapterhead}{\booktermcuestrue}{}{\PackageError{book-links}{Cannot restore starred chapter body cues}{}}
\pretocmd{\@starttoc}{\booktermcuesfalse}{}{\PackageError{book-links}{Cannot suppress contents cues}{}}
\apptocmd{\@starttoc}{\booktermcuestrue}{}{\PackageError{book-links}{Cannot restore contents body cues}{}}
\makeatother

% 手動の式番号は表示を保ち、参照先を式ごとのIDで区別する。
% 番号から最初の参照へ、参照から元の式へ戻れる。
\newcommand{\booktag}[2]{%
  \tag*{\(\cdots\) (\protect\hyperlink{eqcite:#1}{#2})}\label{eqtag:#1}%
}
\DeclareRobustCommand{\bookeqref}[2]{\hyperref[eqtag:#1]{#2}}
\DeclareRobustCommand{\bookeqbackref}[2]{%
  \hypertarget{eqcite:#1}{\bookeqref{#1}{#2}}%
}
\pdfstringdefDisableCommands{%
  \def\bookterm#1#2{#2}%
  \def\booklinkcue#1{#1}%
  \def\booknavref#1#2{#2}%
  \def\bookeqref#1#2{#2}%
  \def\bookeqbackref#1#2{#2}%
}

% 【例題】から【解答】へ、【解答】・【方針】・【解法】から【例題】へ。
% 第1引数は問題群ごとに固有のID。番号のない例題にも使える。
\newcommand{\bookblocklink}[4]{%
  \hypertarget{book:#1:#2}{\hyperlink{book:#1:#3}{\textbf{#4}}}%
}

% enumerate の format に指定し、既存の番号・インデントを保つ。
% 問題番号は解答へ、それ以外の番号は元の問題へリンクする。
% 各問題群の problem / answer / hint / solution で番号をそろえる。
\newcommand{\bookitemlink}[4]{%
  \hypertarget{book:#1:#2:\arabic{enumi}}{%
    \hyperlink{book:#1:#3:\arabic{enumi}}{#4}%
  }%
}
% 既存の問題番号を本文・学習経路から参照する。
% 第1引数は問題群ID、第2引数は番号、第3引数は表示する文字。
\newcommand{\bookproblemref}[3]{%
  \hyperlink{book:#1:problem:#2}{#3}%
}

% 用語の初出表示・巻末の参照案内・式番号のリンクは共通設定で扱う。

\hypersetup{
  pdftitle={漸化式の新解釈},
  pdfauthor={MAS256},
  pdfsubject={線形代数と複素平面から読み直す漸化式}
}

% =========================================================
% 章・節・項の相互参照
% =========================================================

% \headingref{ラベル} : 参照先の番号とタイトルを自動表示する。
% chapter       -> 第n章「タイトル」
% section       -> 第n.m節「タイトル」
% subsection    -> 第n.m.l項「タイトル」
% subsubsection -> 第n.m.l.k項「タイトル」
% 同じ章の直前・直後のsectionだけは「前節」「次節」とする。
% 現在位置がsubsection以下でも、その位置を含むsectionを基準にする。
% \headingref*{ラベル} : 隣接していても番号とタイトルを表示する。
% 既存の\chapterref・\sectionrefも同じ規則を使う（*付きも可）。
% 見出し直後に\labelを置き、番号・タイトルの更新後は再コンパイルする。
% \refは従来どおり番号だけを返す。式・例題の参照にもそのまま使える。
\ExplSyntaxOn
\seq_new:N \l_bookref_number_seq
\tl_new:N \l_bookref_number_tl
\tl_new:N \l_bookref_text_tl

\cs_new_protected:Npn \bookref_full:n #1
  {
    \tl_set:Nn \l_bookref_text_tl { 第\ref*{#1} }
    \int_case:nnF { \seq_count:N \l_bookref_number_seq }
      {
        {1}{ \tl_put_right:Nn \l_bookref_text_tl {章} }
        {2}{ \tl_put_right:Nn \l_bookref_text_tl {節} }
      }
      { \tl_put_right:Nn \l_bookref_text_tl {項} }
    \tl_put_right:Nn \l_bookref_text_tl {「\nameref*{#1}」}
  }

\cs_new_protected:Npn \bookref_format:n #1
  {
    \tl_set:Nx \l_bookref_number_tl { \getrefnumber{#1} }
    \seq_set_split:NnV \l_bookref_number_seq {.} \l_bookref_number_tl
    \bookref_full:n {#1}
    \int_compare:nNnT { \seq_count:N \l_bookref_number_seq } = {2}
      {
        \str_if_eq:eeT
          { \seq_item:Nn \l_bookref_number_seq {1} } { \thechapter }
          {
            \int_compare:nNnT { \value{section} } > {0}
              {
                \int_compare:nNnT
                  { \seq_item:Nn \l_bookref_number_seq {2} }
                  = { \value{section} - 1 }
                  { \tl_set:Nn \l_bookref_text_tl {前節} }
                \int_compare:nNnT
                  { \seq_item:Nn \l_bookref_number_seq {2} }
                  = { \value{section} + 1 }
                  { \tl_set:Nn \l_bookref_text_tl {次節} }
              }
          }
      }
  }

\NewDocumentCommand{\headingref}{s m}
  {
    \group_begin:
    \booktermcuesfalse
    \cs_if_exist:cTF { r@#2 }
      {
        \bookref_format:n {#2}
        \IfBooleanT{#1}{ \bookref_full:n {#2} }
        \hyperref[#2]{\tl_use:N \l_bookref_text_tl}
      }
      { \ref{#2}「\nameref{#2}」 }
    \group_end:
  }
\NewDocumentCommand{\chapterref}{s m}
  { \IfBooleanTF{#1}{\headingref*{#2}}{\headingref{#2}} }
\NewDocumentCommand{\sectionref}{s m}
  { \IfBooleanTF{#1}{\headingref*{#2}}{\headingref{#2}} }
\ExplSyntaxOff

% 紙面の幅で改行位置だけを切り替える。第1引数は本編、第2引数はスマホ版。
\newcommand{\BookMathLayout}[2]{\ifdefined\BookMobile#2\else#1\fi}
\title{漸化式の新解釈}

\ifdefined\BookMobile
  % 同じ原稿をスマホ幅で組版する。入口は scripts/build-mobile.py が生成する。
  \setlength{\mathindent}{1zw}
  \setlength{\fullwidth}{\textwidth}
  \setlength{\columnseprule}{0pt}

  % 文字サイズは保ち、縦の余白を元の約6割にする。
  % 本文10pt・行送り16.8ptから、10pt + 0.6(16.8pt - 10pt) = 14.08ptへ。
  % 行送りそのものを0.6倍すると文字が重なるため、余白だけを詰める。
  \linespread{0.88}
  \newcommand{\BookMobileTightenSkip}[1]{\multiply#1 by 3\divide#1 by 5\relax}
  \BookMobileTightenSkip{\parskip}
  % 脚注のあるページ末尾では、段落の余白をさらに1ptまで詰められるようにする。
  \addtolength{\parskip}{0pt minus 1pt}
  \BookMobileTightenSkip{\smallskipamount}
  \BookMobileTightenSkip{\medskipamount}
  \BookMobileTightenSkip{\bigskipamount}
  \BookMobileTightenSkip{\jot}
  \BookMobileTightenSkip{\skip\footins}
  \BookMobileTightenSkip{\normallineskip}
  \setlength{\lineskip}{\normallineskip}

  % サイズ変更・数式開始時の再設定にも対応する。各コマンドが設定した余白を縮める。
  \newcommand{\BookMobileDisplaySpacing}{%
    \BookMobileTightenSkip{\abovedisplayskip}%
    \BookMobileTightenSkip{\belowdisplayskip}%
    \BookMobileTightenSkip{\abovedisplayshortskip}%
    \BookMobileTightenSkip{\belowdisplayshortskip}}
  \AddToHook{cmd/normalsize/after}{\BookMobileDisplaySpacing}
  \AddToHook{cmd/small/after}{\BookMobileDisplaySpacing}
  \AddToHook{cmd/footnotesize/after}{\BookMobileDisplaySpacing}
  \BookMobileDisplaySpacing

  \renewtcolorbox{bookpoint}{unbreakable,sharp corners,
    colback=black!3,colframe=black!55,boxrule=0pt,leftrule=1.2pt,
    left=5pt,right=3pt,top=1.8pt,bottom=1.8pt,
    before skip=3.6pt,after skip=3.6pt}
  \renewcommand*{\chapterheadstartvskip}{\vspace*{1.38\baselineskip}}
  \renewcommand*{\chapterheadendvskip}{\vspace{1.02\baselineskip}}

  \makeatletter
  % 箇条書き・引用の局所指定も、実際のリスト開始時にまとめて調整する。
  \patchcmd{\list}{\@trivlist}{%
    \BookMobileTightenSkip{\topsep}%
    \BookMobileTightenSkip{\partopsep}%
    \BookMobileTightenSkip{\parsep}%
    \BookMobileTightenSkip{\itemsep}%
    \@trivlist
  }{}{\PackageError{book-mobile}{List spacing patch failed}{Check the LaTeX list definition.}}
  % 見出しの前後だけを詰め、段落見出しの後の横方向の空白は保つ。
  \let\BookMobileStartSection\@startsection
  \newskip\BookMobileHeadingBefore
  \newskip\BookMobileHeadingAfter
  \renewcommand{\@startsection}[6]{%
    \BookMobileHeadingBefore=#4\relax
    \BookMobileHeadingAfter=#5\relax
    \BookMobileTightenSkip{\BookMobileHeadingBefore}%
    \ifdim\BookMobileHeadingAfter>\z@
      \BookMobileTightenSkip{\BookMobileHeadingAfter}%
    \fi
    \BookMobileStartSection{#1}{#2}{#3}%
      {\BookMobileHeadingBefore}{\BookMobileHeadingAfter}{#6}}
  \renewcommand*{\l@subsection}{\@dottedtocline{2}{2zw}{3.5zw}}
  \renewcommand*{\l@subsubsection}{\@dottedtocline{3}{3zw}{4.5zw}}
  \pagestyle{headings}
  \def\@evenhead{\underline{\hbox to\textwidth{\autoxspacing\textbf{\thepage}\hfil
    \parbox[b]{\dimexpr\textwidth-2zw\relax}{\raggedleft\baselineskip=10pt\leftmark}}}}
  \def\@oddhead{\underline{\hbox to\textwidth{\autoxspacing
    \parbox[b]{\dimexpr\textwidth-2zw\relax}{\raggedright\baselineskip=10pt\rightmark}\hfil\textbf{\thepage}}}}
  \makeatother
  % 図だけは、紙幅を超えた場合に描画全体を収める。
  \newsavebox{\BookMobileGraphBox}
  \renewcommand{\BookGraph}[2]{%
    \par\smallskip\noindent
    \begin{minipage}{\linewidth}\centering
      \sbox{\BookMobileGraphBox}{%
        \begin{tikzpicture}[scale=0.8,>=Stealth,line cap=round,line join=round,
          every node/.style={font=\small}]#2\end{tikzpicture}}%
      \ifdim\wd\BookMobileGraphBox>\linewidth
        \resizebox{\linewidth}{!}{\usebox{\BookMobileGraphBox}}%
      \else\usebox{\BookMobileGraphBox}\fi
      \BookFigureCaption{#1}%
    \end{minipage}\par\smallskip
  }
\fi

% 両版のログから、はみ出した式の入力ファイルを特定する。
\AddToHook{file/before}{\typeout{BOOK-BEGIN:\CurrentFilePath\CurrentFile}}
\AddToHook{file/after}{\typeout{BOOK-END:\CurrentFilePath\CurrentFile}}


